%% =====================================================================
%%  WHSS Google Day @ IISc Poster
%% =====================================================================

\documentclass[a0paper,portrait,fleqn]{beamer}
\usepackage[orientation=portrait,size=a0,scale=1.0,debug=false]{beamerposter}
\usepackage[T1]{fontenc}
\usepackage[utf8]{inputenc}
\usepackage[english]{babel}
\usepackage{helvet}
\renewcommand{\familydefault}{\sfdefault}
\usepackage{graphicx}
\usepackage{tikz}
\usetikzlibrary{arrows.meta,positioning,shapes.geometric,fit,shadows}
\usepackage{pgfplots}
\pgfplotsset{compat=1.18}
\usepackage{booktabs}
\usepackage{array}
\usepackage{colortbl}
\usepackage{xcolor}
\usepackage{enumitem}
\usepackage{qrcode}
\usepackage{microtype}
\usepackage[most]{tcolorbox}
\usepackage{hyperref}
\hypersetup{colorlinks=true, urlcolor=blue!60!black, linkcolor=blue!60!black}
\usepackage{calc}

% Colors
\definecolor{gblue}{HTML}{4285F4}
\definecolor{gred}{HTML}{EA4335}
\definecolor{gyellow}{HTML}{FBBC05}
\definecolor{ggreen}{HTML}{34A853}
\definecolor{darktext}{HTML}{1A1A1A}
\definecolor{lightgray}{HTML}{F5F5F5}
\definecolor{midgray}{HTML}{5F6368}

\setbeamercolor{background canvas}{bg=white}
\setbeamercolor{normal text}{fg=darktext}
\setbeamertemplate{navigation symbols}{}

\newcommand{\postertitlesize}{\fontsize{64}{76}\selectfont}   
\newcommand{\posterauthorsize}{\fontsize{34}{40}\selectfont}
\newcommand{\tracksize}{\fontsize{26}{30}\selectfont}
\newcommand{\sectionheadsize}{\fontsize{40}{46}\selectfont}   
\newcommand{\subheadsize}{\fontsize{30}{35}\selectfont}       
\newcommand{\bodysize}{\fontsize{24}{31}\selectfont}          
\newcommand{\captionsize}{\fontsize{18}{22}\selectfont}       

\setbeamerfont{block title}{size=\Large,series=\bfseries}
\setbeamerfont{block body}{size=\normalsize}
\renewcommand{\normalsize}{\bodysize}

\renewenvironment{block}[1]{%
  \begin{tcolorbox}[breakable, enhanced, 
    colback=white, colframe=gblue,
    colbacktitle=gblue, coltitle=white,
    boxrule=2.5pt, arc=12pt, auto outer arc,
    titlerule=0pt,
    drop fuzzy shadow=black!25,
    fonttitle=\sectionheadsize\bfseries,
    title={#1}, top=18pt, bottom=18pt, left=20pt, right=20pt,
    toptitle=14pt, bottomtitle=14pt, lefttitle=18pt,
    before skip=24pt, after skip=14pt]
  \color{darktext}\bodysize
}{%
  \end{tcolorbox}
}

\setlength{\parskip}{10pt}
\setlist[itemize]{leftmargin=32pt, itemsep=8pt, label=\textcolor{gblue}{\textbullet}}

\begin{document}
\begin{frame}[t]
\vspace{20pt}

\begin{beamercolorbox}[wd=\linewidth, sep=10pt]{header}
\begin{minipage}[c]{0.14\linewidth}
  \centering
  \includegraphics[width=1.3\linewidth]{IISc_Master_Seal_Transparent.png}
\end{minipage}
\hfill
\begin{minipage}[c]{0.68\linewidth}
  \begin{tcolorbox}[enhanced, colback=gblue!5!white, colframe=gblue!30!white, 
    boxrule=2pt, arc=15pt, auto outer arc,
    top=24pt, bottom=24pt, left=20pt, right=20pt,
    drop fuzzy shadow=black!10]
    \centering
    \color{gblue}\postertitlesize\textbf{Warped Hybrid Slice Sampling (WHSS):}\\[10pt]
    \color{darktext}\textbf{Scaling Constrained Posteriors via Geometric Preconditioning}\\[20pt]
    \posterauthorsize
    \underline{Dharmik Patel}$^{1}$, Prof.\ Murugesan Venkatapathi (Advisor)$^{1}$\\[6pt]
    \bodysize\color{midgray}
    $^{1}$Computation, Statistics and Physics Lab (CSPL), Dept.\ of Computational and Data Sciences (CDS), IISc Bengaluru\\[14pt]
    \tracksize\color{gblue}\textbf{Research Track 1: Core AI/ML \& Trustworthy Systems (Theory \& Algorithms)}
  \end{tcolorbox}
\end{minipage}
\hfill
\begin{minipage}[c]{0.14\linewidth}
  \centering
  \includegraphics[width=\linewidth]{CDS_Logo_Wide.png}
\end{minipage}
\end{beamercolorbox}

\vspace{4pt}
{\color{gblue}\rule{\linewidth}{4pt}}
\vspace{8pt}

\newlength{\pgutter}
\setlength{\pgutter}{28pt}
\newlength{\pcolwidth}
\setlength{\pcolwidth}{(\linewidth-2\pgutter)/3}

% --- COLUMN 1 ---
\begin{minipage}[t]{\pcolwidth}

\begin{block}{Motivation}
Sampling is the backbone of statistical algorithms across diverse fields, including modern \textbf{Generative AI}. However, many real-world applications require sampling from strictly \textbf{non-differentiable (black-box) target densities} bounded by \textbf{highly skewed geometric constraints} ($Aw \le b$). Navigating these opaque, rigid polytopes presents a massive computational bottleneck. To bridge this gap, we propose WHSS and demonstrate its \textbf{empirical dominance on real-world problems} like Institutional CVaR Portfolios and genome-scale Metabolic Flux Analysis.
\end{block}

\begin{block}{Limitations of Prior SOTA}
\begin{itemize}
  \item \textbf{Hit-and-Run / Dikin Walk:} Suffer exponential efficiency decay in high dimensions ($\ge 100D$) due to isotropic step proposals and $\mathcal{O}(D^3)$ computational bottlenecks.
  \item \textbf{Gradient-Based (NUTS/HMC):} Require strictly differentiable target densities (failing on black-box models) and struggle to navigate hard inequality boundaries without extreme tuning.
  \item \textbf{Emcee (Ensemble):} Affine-invariance mathematically breaks down in spaces with overlapping sharp corners, leading to permanent spatial trapping and pseudo-convergence.
\end{itemize}
\end{block}

\begin{block}{Core Hypothesis}
We hypothesize that by completely decoupling geometric boundary discovery (via vMF ray-casting) from density exploration (via parallel Hybrid Slice Sampling), we can construct a global affine warp matrix $L$ that insulates the Markov chain from dimensional collapse, maintaining a constant mixing rate regardless of ambient dimensionality.
\end{block}

\begin{block}{Key Contributions}
\begin{itemize}
  \item \textbf{Amortized Scalability:} Achieves dimension-independent MCMC mixing efficiency (up to 400D) at an amortized $\mathcal{O}(D^2)$ computational cost.
  \item \textbf{Robust Hybrid Exploration:} Bypasses the ``ensemble collapse'' trap of prior SOTAs (Emcee) by coupling exact boundary clipping with novel \textbf{Skeleton proposals} to guarantee global mobility—empirically proven by superior Mean Squared Jump Distance (MSJD) and rapid ACF decorrelation.
  \item \textbf{Production Validation:} Demonstrated empirical dominance over SOTAs on 30D Finance CVaR risk targets and 24D Biological Transcriptomics integration.
\end{itemize}
\end{block}

\begin{block}{Open Source Artifacts}
\begin{minipage}[c]{0.3\linewidth}
\centering
\vspace{4pt}
\qrcode[height=2.8cm]{https://github.com/user/whss}
\vspace{4pt}
\end{minipage}%
\begin{minipage}[c]{0.65\linewidth}
\textbf{\color{gblue}Scan for Code \& Experiments}\\[6pt]
\small Reproduce all benchmarks, MCMC traces, and results.
\end{minipage}
\end{block}

\end{minipage}%
\hspace{\pgutter}%
% --- COLUMN 2 ---
\begin{minipage}[t]{\pcolwidth}

\begin{block}{System Architecture: The WHSS Pipeline}
\vspace{10pt}
\begin{center}
\begin{tikzpicture}[
    node distance=1.1cm and 1.1cm,
    box/.style={draw=gblue, line width=2.5pt, rounded corners=8pt, fill=white,
                minimum width=7.5cm, minimum height=1.8cm, align=center, font=\bodysize, drop fuzzy shadow=black!20},
    arrow/.style={-{Latex[length=12pt]}, line width=2.5pt, gblue}
  ]
  \node[box] (input) {Input Space\\$Aw \le b$ \& Density $f(x)$};
  \node[box, below=of input, fill=gred!15] (phase0) {Phase 0: Deterministic Axis Scouts\\(Structural Boundary Discovery)};
  \node[box, below=of phase0, fill=gyellow!25] (phase1) {Phase 1: vMF Ray-Casting\\(Statistical Density Discovery)};
  \node[box, below=of phase1, fill=gyellow!25] (phase2) {Phase 2: Adaptive Preconditioning\\(Convex Shrinkage $L$-Matrix)};
  \node[box, below=of phase2, fill=ggreen!20] (phase3) {Phase 3: Hybrid Slice Sampling\\(Warped, Skeleton \& Coord Rays)};
  \node[box, below=of phase3] (output) {Output:\\Ergodic MCMC Samples};
  \draw[arrow] (input) -- (phase0);
  \draw[arrow] (phase0) -- (phase1);
  \draw[arrow] (phase1) -- (phase2);
  \draw[arrow] (phase2) -- (phase3);
  \draw[arrow] (phase3) -- (output);
\end{tikzpicture}
\end{center}
\vspace{10pt}
\end{block}

\begin{block}{Algorithmic Novelty}
\begin{itemize}
  \item \textbf{Geometric Preconditioning:} WHSS applies adaptive convex shrinkage to anchor covariance, learning a robust warp matrix $L$ that transforms both density and constraints into an isotropic space:
  \vspace{-10pt}
  \[ x = L^{-1}(w - \mu) \implies f(Lx + \mu) \text{ s.t. } (AL)x \le b - A\mu \]
  \vspace{-15pt}
  \item \textbf{Hybrid Proposal Engine (Skeleton Rays):} WHSS dynamically mixes Affine Warped rays, Coordinate rays, and novel \textbf{Skeleton rays}. By calculating jump vectors between static Phase 1 boundary anchors (rather than active MCMC states), WHSS permanently bypasses the correlation penalties and ensemble collapse of standard Differential Evolution.
  \vspace{-10pt}
  \item \textbf{Parallel Adaptive Refinement (PATT):} 10 independent chains update the global $L$-matrix using Welford's recursion, freezing strictly at tick $k_{\max}=10$ to satisfy detailed balance and ergodicity.
\end{itemize}
\end{block}

\begin{block}{Algorithmic Limitations}
\begin{itemize}
  \item \textbf{Needle Trap (Geometry):} WHSS is vulnerable on purely uniform distributions inside extremely long, thin polytopes. Isotropic vMF warm-up rays collide with side-walls, generating a collapsed $L$-matrix and trapping the chain. CHRR remains the standard for such cases.
  \item \textbf{Warm-Up Cost (Speed):} WHSS is inherently slower than single-chain baselines (e.g., 13s vs 0.6s for Emcee in 30D). The $\mathcal{O}(D^3)$ Phase 0/1 warm-up is a one-time investment amortized over long chains, but it is a real cost for short sampling budgets.
\end{itemize}
\end{block}

\end{minipage}%
\hspace{\pgutter}%
% --- COLUMN 3 ---
\begin{minipage}[t]{\pcolwidth}

\begin{block}{Empirical Validation 1: Dimensional Scaling}
\begin{center}
\begin{tikzpicture}
\begin{axis}[
    width=0.95\linewidth, height=0.62\linewidth,
    xlabel={Dimensionality ($D$)},
    ylabel={Efficiency (ESS / 1k NFE)},
    label style={font=\captionsize, font=\bfseries},
    tick label style={font=\captionsize},
    legend style={font=\captionsize, at={(0.95,0.95)}, anchor=north east,
                  legend columns=1, draw=midgray, fill=white, rounded corners=2pt},
    legend cell align={left},
    ymin=0.1, ymax=100,
    xmin=10, xmax=400,
    xmode=log, ymode=log,
    xtick={10, 50, 100, 200, 400},
    xticklabels={10, 50, 100, 200, 400},
    grid=both, grid style={dashed, gray!40},
    minor grid style={dotted, gray!20},
    axis x line*=bottom,
    axis y line*=left,
    enlarge x limits=0.05,
    line width=1.8pt,
  ]
  \addplot[gblue, mark=*, mark size=3.5pt, mark options={fill=white, draw=gblue, line width=1.5pt}] coordinates
    {(10,36.44) (50,7.43) (100,5.51) (200,5.32) (400,4.40)};
  \addplot[gyellow, mark=square*, mark size=3.5pt, mark options={fill=white, draw=gyellow, line width=1.5pt}] coordinates
    {(10,30.65) (50,1.99) (100,0.83) (200,0.58) (400,0.48)};
  \addplot[gred, mark=triangle*, mark size=4.0pt, mark options={fill=white, draw=gred, line width=1.5pt}] coordinates
    {(10,17.09) (50,1.05) (100,0.62) (200,0.84) (400,0.44)};
  \legend{{WHSS (Ours, Post-Warm-Up)}, {Hit-and-Run (HRSS)}, {Dikin Walk}}
\end{axis}
\end{tikzpicture}
\end{center}
\vspace{10pt}
\textbf{\color{gblue}Result (Dimensional Stability):} WHSS maintains a steady efficiency floor of $\sim$4.40 ESS up to 400D, whereas prior SOTAs suffer exponential efficiency decay.

\vspace{6pt}
\textbf{\color{darktext}Methodological Note:} Tested on a highly skewed Gaussian ($\kappa=1000$) inside an isotropic hypercube. This formally isolates and proves WHSS scales efficiently against \textbf{target-skew}, assuming the underlying constraint geometry is well-behaved.
\end{block}

\begin{block}{Real-World Impact: Finance \& Bioinformatics}
\textbf{\color{darktext}1. Finance:} 30-Asset Portfolio with Sector Caps \& Non-Differentiable Risk.
\begin{center}
\renewcommand{\arraystretch}{1.2}
\begin{tabular}{lcc}
\toprule
\rowcolor{gblue!15}
\textbf{Method} & \textbf{ESS/1k NFE} & \textbf{Domain Coverage} \\
\midrule
Emcee (Ensemble) & 11.15$^\dagger$ & 79.3\% (\textcolor{gred}{Trapped}) \\
HRSS      & 2.80 & 100.0\% \\
Dikin Walk & 2.91 & 99.7\% \\
\textbf{WHSS (Ours)} & \textbf{3.51} & \textbf{100.0\% (\textcolor{ggreen}{Global})} \\
\bottomrule
\end{tabular}
\end{center}
\textbf{\color{gred}Failure Mode ($^\dagger$):} Emcee's affine-invariance breaks on sharp intersecting constraints.

\vspace{16pt}
\textbf{\color{darktext}2. Bioinformatics:} Transcriptomics (Gaussian penalty) in 24D.
\begin{center}
\renewcommand{\arraystretch}{1.2}
\begin{tabular}{lcc}
\toprule
\rowcolor{gblue!15}
\textbf{Method} & \textbf{Weight Ret.} & \textbf{Min ESS} \\
\midrule
CHRR (Uniform + SIR) & 0.01\% & 1.0 (Degenerate) \\
\textbf{WHSS (Ours)} & \textbf{100.0\%} & \textbf{102.8} \\
\bottomrule
\end{tabular}
\end{center}
\textbf{\color{gred}Failure Mode:} SIR (Sampling Importance Resampling) completely degenerates on skewed targets.

\vspace{16pt}
\textbf{\color{gblue}Core Observation:} Across both financial and biological domains, prior algorithms either get trapped in sharp geometric corners (Emcee) or catastrophically collapse when approximating skewed target densities (CHRR+SIR). WHSS maps these spaces natively without relaxation.
\end{block}

\end{minipage}

\end{frame}
\end{document}
